\documentclass[11pt]{article}
\usepackage[margin=1in]{geometry}
\usepackage{amsmath,amssymb,amsthm,mathtools}
\usepackage{booktabs}
\usepackage[expansion=false]{microtype}
\usepackage{xcolor}
\usepackage[colorlinks=true,linkcolor=blue!50!black,citecolor=blue!50!black,urlcolor=blue!50!black]{hyperref}
\setlength{\emergencystretch}{3em}

\newtheorem{theorem}{Theorem}[section]
\newtheorem{proposition}[theorem]{Proposition}
\newtheorem{lemma}[theorem]{Lemma}
\newtheorem{corollary}[theorem]{Corollary}
\theoremstyle{definition}
\newtheorem{definition}[theorem]{Definition}
\theoremstyle{remark}
\newtheorem{remark}[theorem]{Remark}

\newcommand{\E}{\mathbb E}
\newcommand{\Var}{\operatorname{Var}}
\newcommand{\Tr}{\operatorname{Tr}}
\newcommand{\Res}{\operatorname{Res}}

\title{Two objects\\[4pt]
\large a random walk on the tree of activation codes, an isotropic process on its Cantor set,\\
and activation means as Dixmier traces}
\author{Working note VIII for the continuation-algebra programme\\
\small self-reviewed draft, not independently verified}
\date{3 October 2026}

\begin{document}
\maketitle

\begin{abstract}
The proposal has two separate objects.
\begin{itemize}
\item \emph{Object 1}: nodes are binary strings (activation codes), and arrows are shift operators with self-loops, so the structure is a groupoid with a convolution algebra.
\item \emph{Object 2}: the tree of codes, carrying values defined from the weights, whose Cantor-set distance (last point of contact) should translate into a process, and a run length, for estimating activation means. The Dixmier trace and the zeta function are expected to play a role.
\end{itemize}
Using Bendikov--Grigor'yan--Pittet--Woess (BGPW) on isotropic Markov semigroups on ultrametric spaces, and the path correspondences of quantum graphs (Khatua--Roy), we make this exact.

\textbf{(1) Object 1 generates object 2.} The de Bruijn code graph, with self-loops at $0^n$ and $1^n$, is a quantum graph whose path correspondences are the levels of the code tree. Its convolution algebra is the Cuntz algebra $\mathcal O_2$: nodes are projections $P_v=S_vS_v^*$ indexed by binary strings, and arrows are the shift isometries $S_0,S_1$.

\textbf{(2) A walk on object 1 samples object 2.} A nearest-neighbour walk on the code tree that backtracks with depth-dependent probability, and otherwise steps forward with the conditional code law, exits with \emph{exactly} the code law $\mu$. Verified: $\chi^2/\mathrm{dof}=1.02$ over $64$ cylinders.

\textbf{(3) Duality (BGPW, Kigami).} The induced process on the Cantor set is isotropic. Its jump kernel depends only on the last point of contact $x\wedge y$, and its spectrum is $\{1/G(v,o)\}$, given by the Green function of the walk on object~1.

\textbf{(4) Weights $\to\kappa\to$ zeta and Dixmier.} With the ultrametric $d(x,y)=\mu(x\wedge y)^{1/s}$, built from the weight-induced code law, the isotropic Laplacian $L_s$ has spectral zeta function $\zeta_s(t)=\Tr L_s^{-t}$. Its simple pole is at $t=s$, with \emph{residue $s/h$, where $h$ is the entropy rate of the code process}. Its Dixmier trace is integration against $\mu$, so \emph{activation means are Dixmier traces}. The residue and the subtree ratios are checked numerically, the latter to $3\times10^{-4}$.

\textbf{(5) Run length.} The process runs for time $T(\varepsilon)=2\varepsilon^{-2}\sum_vG(v,o)\,e_v$, a Green-function-weighted energy.

\textbf{(6) The verdict.} A ReLU network on Gaussian input has no internal randomness, so every fresh code costs a fresh forward pass. Code cells chosen from the weights carry at most about $10\%$ of the variance at any layer: stratification gains are $\le1.10\times$. The two-phase criterion $R^2_\ell>\ell/L$ fails everywhere. The framework gives an exact NCG formula for the means, but not a faster estimator for random He networks.
\end{abstract}

\section{Object 1: the code graph and its algebra}

Read the network's neuron signs sequentially, neuron by neuron and layer by layer. Each input gives a binary string, a path in the binary \emph{code tree} $T$, of depth $nL$. The window of the last $n$ signs moves by the de Bruijn shift $s_1\cdots s_n\mapsto s_2\cdots s_nb$. After $n$ shifts it holds the next layer's code; the position within the layer is a phase modulo $n$, and this is ``shift modulo cyclicness''. The constant codes $0^n$ (a dead layer) and $1^n$ (a fully active layer) carry self-loops.

\begin{proposition}[object 1 builds object 2]\label{prop:o1}
Take the de Bruijn graph on $\{0,1\}^n$, viewed as a classical quantum graph $(C(\{0,1\}^n),\psi,A)$ with uniform $\psi$ and $0/1$ adjacency $A$. Its $k$-path correspondences are $E^k=C(\{\text{paths of length }k\})\cong C(\{0,1\}^{n+k})$ (Khatua--Roy, \S4.1), i.e.\ the levels of the binary tree. The infinite path space is the Cantor set $\{0,1\}^{\mathbb N}$ with the one-sided shift. The associated Cuntz--Krieger algebra is the Cuntz algebra $\mathcal O_2$, because the edge shift is the full $2$-shift. In it, nodes are the projections $P_v=S_vS_v^*$ indexed by binary strings $v$, and arrows are the isometries $S_0,S_1$. A nonnegative function on arrows of the shift groupoid that sums to one along the arrows leaving each point, $\sum_b p(b\mid\cdot)=1$, acts by the Ruelle transfer operator $f\mapsto\sum_bp(b\mid\cdot)f(\cdot b)$: convolution is one Markov step.
\end{proposition}

Paths in object~1 are vertices of object~2. Object~1's Markov steps are transfer operators, which are the moves of a walk on the tree.

\section{Object 1's walk samples object 2}

\begin{theorem}[exit measure]\label{thm:exit}
Let $\mu$ be any probability on the code Cantor set with $\mu(v)>0$ for every vertex. Let $0<q_k<1/2$ for $k\ge1$ and $q_0=0$. The nearest-neighbour walk on $T$ steps back to the parent with probability $q_{|v|}$, and otherwise moves to the child $vb$ with probability $(1-q_{|v|})\,\mu(vb)/\mu(v)$. It is transient, and its exit measure from the root is exactly $\mu$.
\end{theorem}

\begin{proof}
The depth process is a birth--death chain that does not depend on $\mu$ and drifts outward. So the probability $\epsilon_{k}$ of never returning to the parent, starting from depth $k$, depends only on $k$. The walk exits through the cylinder $[v]$ iff its last departure from $v$ goes to some child and never returns. That child is $vb$ with probability proportional to $(1-q_{|v|})\frac{\mu(vb)}{\mu(v)}\epsilon_{|v|+1}$, i.e.\ to $\mu(vb)/\mu(v)$. Induct on depth.
\end{proof}

\textbf{Check} (\texttt{code/tree\_walk\_exit.py}). On a random binary $\mu$ of depth $6$, with $q_k=0.3+0.1\sin k$ and $4\times10^5$ walks, the maximum error is $9\times10^{-4}$ and $\chi^2/\mathrm{dof}=1.02$ over $64$ cylinders.

\section{Duality: the process on object 2}

\begin{theorem}[BGPW Thms.~3.8, 7.1, 7.5; Kigami]\label{thm:dual}
The boundary process induced by a transient nearest-neighbour walk on $T$, with reference measure its exit measure $\mu$, is an \emph{isotropic} jump process on $\partial T$: its jump kernel is $J(x,y)=j(x\wedge y)$, a function of the last point of contact only. Its Laplacian has pure point spectrum $\{1/G(v,o)\}\cup\{0\}$, given by the Green function of the walk. The eigenfunctions are $f_C=\mathbf 1_C/\mu(C)-\mathbf 1_B/\mu(B)$ for $C$ a child of $B$, with eigenvalue determined by $B$. Conversely, every isotropic process, for any reference measure and any ultrametric, arises in this way from a unique walk on $T$, up to a linear time change.
\end{theorem}

The eigenspaces $H_B=\mathrm{span}\{f_C:C\prec B\}$ are the detail spaces of note~VII. That note's spectral formula for last-point-of-contact kernels is the finite-tree, discrete-time case of BGPW's Theorem~3.8.

\section{Weights, zeta, Dixmier: activation means as Dixmier traces}

The weights determine the code law $\mu$ under Gaussian input. Given $s>0$, take the ultrametric
\[d_s(x,y)=\mu(x\wedge y)^{1/s}\qquad(\text{Michon weights }\kappa=\mu^{1/s}).\]
By Theorem~\ref{thm:dual}, the BGPW isotropic Laplacian $L_s$ of $(d_s,\mu)$ has eigenvalue $\mu(B)^{-1/s}$ on $H_B$, up to the normalization of the distance distribution, and is realized by a unique walk on object~1.

\begin{theorem}[zeta, entropy, Dixmier]\label{thm:zeta}
Suppose the code law is a finite-memory (de Bruijn) Markov process with positive entropy rate $h$. The general sequential code law is not of this form; the theorem applies to its Markov approximation, or to the infinite-depth limit when that limit is Gibbs. On the binary tree, let
\[\zeta_s(t)=\Tr\big(L_s^{-t}\big|_{1^\perp}\big)=\sum_B\mu(B)^{t/s}.\]
\begin{enumerate}
\item[(i)] $\zeta_s$ converges for $t>s$ and has a simple pole at $t=s$, with $\Res_{t=s}\zeta_s=s/h$.
\item[(ii)] For a vertex $v$, the subtree series $\zeta_{s,v}(t)=\sum_{B\subseteq[v]}\mu(B)^{t/s}$ satisfies $\Res\zeta_{s,v}/\Res\zeta_s=\mu(v)$.
\item[(iii)] Hence, for continuous $f$ on the code Cantor set,
\[\int f\,d\mu=\lim_{t\downarrow s}\frac{\Tr\big(M_fL_s^{-t}\big)}{\Tr\big(L_s^{-t}\big)}=\frac{\Tr_\omega\big(M_fL_s^{-s}\big)}{\Tr_\omega\big(L_s^{-s}\big)}.\]
The Dixmier trace of $L_s^{-s}$ is integration against the code law.
\end{enumerate}
\end{theorem}

\begin{proof}
With $\sigma=t/s$ and $P$ the transition matrix on de Bruijn states, $\sum_{|B|=k}\mu(B)^\sigma=\pi_\sigma^\top P_\sigma^{k-m}\mathbf 1$, where $P_\sigma$ is the entrywise power. The Perron root $\rho(\sigma)$ is analytic near $\sigma=1$, with $\rho(1)=1$ and
\[\rho'(1)=\pi^\top(P\circ\log P)\mathbf 1=-h.\]
So $\sum_k\rho(\sigma)^k\sim1/(h(\sigma-1))$, which gives (i); the factor $s$ comes from $\sigma-1=(t-s)/s$.

For (ii), the subtree series starts from the state of $v$ with mass $\mu(v)$. At $\sigma=1$ the right Perron vector is $\mathbf 1$, so the residue is $\mu(v)s/h$.

For (iii), $H_B$ is one-dimensional, with a unit eigenvector $\phi_B$ supported in $B$. Then $\Tr(M_fL_s^{-t})=\sum_B\mu(B)^{t/s}\langle\phi_B,f\phi_B\rangle$, and $\langle\phi_B,f\phi_B\rangle$ is a weighted average of $f$ over $B$. As $t\downarrow s$ the deep $B$ dominate, so (ii) and continuity of $f$ give the limit. The Dixmier form is Connes' residue identification, valid for a simple pole.
\end{proof}

\textbf{Check} (\texttt{code/zeta\_dixmier.py}). Memory-$3$ chain, $h=0.5865$ nats, $s=1.7$. As $t\downarrow s$, $(t-s)\zeta_s(t)=2.697,2.816,2.857,2.878$, against $s/h=2.899$ with $O(t-s)$ convergence. The subtree residue ratios match $\mu(v)$ to $3\times10^{-4}$.

\begin{corollary}[activation means]
Let $f_j(c)=\E[F_j\mid\text{code}=c]$ be the conditional mean of output $j$ on the tile with full code $c$ (note~V). Then $\E F_j=\int f_j\,d\mu$ is the Dixmier trace of $M_{f_j}$ against $L_s^{-s}$. The zeta residue sees only the entropy rate of the code process. Within-tile variation, $\E\Var(F_j\mid\text{code})$, is invisible to the code tree.
\end{corollary}

\section{Which process, how long, and what it costs}

\begin{proposition}[run length]\label{prop:run}
Run the stationary isotropic process for time $T$ and average $f$. The variance is asymptotically
\[\frac2T\,\langle f,L^{-1}f\rangle=\frac2T\sum_BG(B,o)\,e_B,\qquad e_B=\|P_{H_B}f\|^2.\]
So reaching accuracy $\varepsilon$ takes time $T(\varepsilon)=2\varepsilon^{-2}\sum_BG(B,o)e_B$, the Green-function pairing of the activation's tree energy with object~1's walk.
\end{proposition}

\textbf{Cost.} A jump with last point of contact at depth $k$ needs a fresh code below $k$. For a network whose only randomness is its input, that means a fresh input conditioned on the shared prefix, which costs at least one pass through the later layers plus conditional sampling. By note~VII (positive spectrum), the process is then never better per evaluation than independent inputs. The only lever is stratification by code cells, priced next.

\begin{proposition}[two-phase stratification along the code tree]\label{prop:twophase}
Let the codes at depth $\ell$ explain a fraction $R^2_\ell$ of $\Var F$ through their cells, and cost a fraction $c_\ell$ of a forward pass to compute. Run $n_1$ inputs to depth $\ell$ to estimate cell masses, and finish $n_2\le n_1$ of them. The best variance-times-cost gain over Monte Carlo is
\[G_\ell=\Big(\sqrt{R^2_\ell c_\ell}+\sqrt{(1-R^2_\ell)(1-c_\ell)}\Big)^{-2}\quad\text{if }R^2_\ell>c_\ell,\qquad G_\ell=1\ \text{otherwise}.\]
\end{proposition}

\begin{proof}
The variance is $\approx(1-R^2)\Var F/n_2+R^2\Var F/n_1$, and the cost is $n_1c+n_2(1-c)$. Minimize under $n_2\le n_1$; the constraint binds exactly when $R^2\le c$.
\end{proof}

\textbf{Measurement} (\texttt{code/code\_tree\_energy.py}). Width $256$, $2^{20}$ inputs, eight live outputs. The table gives the stratification gain $\Var F/\E\Var(F\mid\text{cell})=1/(1-R^2)$ when cells are the activation codes of $k$ neurons chosen from the weights. At layer~$1$ the neurons are ranked by linearized sensitivity of the output; in the middle layer, by the linearized map; at layer $L-1$, by $|W_L[j,i]|$.
\begin{center}\small\setlength{\tabcolsep}{4pt}
\begin{tabular}{lccccccccc}
\toprule
& \multicolumn{3}{c}{layer 1} & \multicolumn{3}{c}{middle} & \multicolumn{3}{c}{layer $L-1$}\\
depth & $k{=}4$ & $8$ & $12$ & $4$ & $8$ & $12$ & $4$ & $8$ & $12$\\
\midrule
4 & 1.025 & 1.041 & 1.056 & 1.017 & 1.035 & 1.057 & 1.042 & 1.071 & 1.100\\
8 & 1.006 & 1.012 & 1.020 & 1.020 & 1.040 & 1.058 & 1.043 & 1.073 & 1.101\\
16 & 1.003 & 1.006 & 1.010 & 1.015 & 1.023 & 1.031 & 1.048 & 1.073 & 1.094\\
\bottomrule
\end{tabular}
\end{center}
At most about $9\%$ of the variance is explained, and it is at the last hidden layer, where $c_\ell=(L-1)/L$. The cheap first-layer cells explain about $1\%$ at depth $16$. Hence $G_\ell\approx1$ everywhere.

\section{Verdict}
\begin{itemize}
\item \emph{The construction is exact.}
\begin{itemize}
\item Object~1 (the de Bruijn code graph, with Cuntz algebra $\mathcal O_2$) generates object~2 (the code tree and its Cantor set) through path correspondences.
\item A walk on object~1 samples object~2's code law exactly (Theorem~\ref{thm:exit}).
\item The induced process on object~2 is BGPW-isotropic, with last-point-of-contact jumps and Green-function spectrum.
\item Weights define $\kappa=\mu^{1/s}$, whose zeta function has residue $s/h$ and whose Dixmier trace is the code law, so activation means are Dixmier traces.
\item The run length is a Green-weighted energy.
\end{itemize}
\item \emph{It does not beat Monte Carlo for random He networks.} Randomness enters only at the input, so every fresh code costs a forward pass, and code cells carry little of the variance (above). Together with notes~II and VII, the variance of deep random networks is spread over exponentially many cells and high Hermite degrees.
\item \emph{Where it would pay.} The structure pays when moving deep in the code tree is cheaper than a full pass. That holds for networks with internal randomness (dropout, noisy units, stochastic depth), where backtracking reuses the prefix computation, and for networks whose code law has low entropy rate $h$, which concentrates the Dixmier measure.
\end{itemize}

\appendix
\section{Scripts}
\texttt{code/tree\_walk\_exit.py} (Theorem~\ref{thm:exit}); \texttt{code/zeta\_dixmier.py} (Theorem~\ref{thm:zeta}); \texttt{code/code\_tree\_energy.py} with output (Proposition~\ref{prop:twophase} inputs).
\end{document}
